\section{The six-cube record}\label{S:six}

\subsection{Maximal cliques and their orbits}
$G_6$ (distance $\le3$) has $10{,}752$ maximal cliques in seven orbits of $B_6$ ($|B_6|=46{,}080$); Table~\ref{tab:S4cl} gives size, orbit size, exact value and a maximiser. The orbit decomposition is certified by comparing the orbit size of each representative (computed by counting stabilising group elements) with the size of its invariant class.
\begin{table}[H]\centering\small
\caption{Maximal cliques of $G_6$ by orbit (main Theorem~\mref{thm:m6}).}\label{tab:S4cl}
\begin{tabular}{rrrll}
\toprule
size & orbit & $V$ & maximiser $y$ & survivable\\
\midrule
12 & 192 & $3$ & $(-1,-1,0,1,1,-1)$ & yes\\
10 & 1{,}280 & $\tfrac{5}{2}$ & $(-1,-1/2,1/2,1/2,1,-1)$ & yes\\
8 & 160 & $3$ & $(0,-1,0,0,1,-1)$ & yes\\
8 & 2{,}880 & $2$ & $(0,0,0,0,1,-1)$ & no\\
8 & 3{,}840 & $2$ & $(0,0,0,0,1,-1)$ & no\\
8 & 1{,}920 & $2$ & $(0,0,0,0,1,-1)$ & no\\
8 & 480 & $2$ & $(0,0,0,0,1,-1)$ & no\\
\bottomrule
\end{tabular}

\end{table}

\subsection{Maximal survivable sets}
Table~\ref{tab:S4max} is the certificate record of the three maximal survivable orbits; the representatives, maximisers and the thresholds of all their nonempty subsets follow.
\begin{table}[H]\centering\small
\caption{Six-cube maximal survivable orbits with certificates; the last column counts the nonempty subsets of the representative with $V=\tfrac52$ exactly.}\label{tab:S4max}
\begin{tabular}{rrrrrrrrlr}
\toprule
size & orbit & $V$ & $f(S)$ & $L^\ast$ & cap & antichain & lost after & word period / dip & boundary subsets\\
\midrule
12 & 192 & $3$ & 1 & 1 & 6 & 22 & 0 & 2 / 1 & 0\\
10 & 1{,}280 & $\tfrac{5}{2}$ & 1 & 3 & 12 & 57 & 2 & 64 / 48 & 112\\
8 & 160 & $3$ & 5 & 5 & 6 & 8 & 0 & 8 / 7 & 0\\
\bottomrule
\end{tabular}

\end{table}
\begin{description}[leftmargin=1.5em,style=nextline]
\item[12 cells, orbit 192.] $y=(-1,-1,0,1,1,-1)$. Cells: \texttt{----+-}, \texttt{---+--}, \texttt{---++-}, \texttt{---+++}, \texttt{--+-+-}, \texttt{--++--}, \texttt{--+++-}, \texttt{--++++}, \texttt{-+-++-}, \texttt{-++++-}, \texttt{+--++-}, \texttt{+-+++-}. Nonempty subsets by threshold: $L^\ast=0$: 266, $L^\ast=1$: 3{,}829.
\item[10 cells, orbit 1{,}280.] $y=(-1,-1/2,1/2,1/2,1,-1)$. Cells: \texttt{----+-}, \texttt{---++-}, \texttt{--+-+-}, \texttt{--++--}, \texttt{--+++-}, \texttt{--++++}, \texttt{-+-++-}, \texttt{-++-+-}, \texttt{-++++-}, \texttt{+-+++-}. Nonempty subsets by threshold: $L^\ast=0$: 161, $L^\ast=1$: 750, $L^\ast=3$: 112.
\item[8 cells, orbit 160.] $y=(0,-1,0,0,1,-1)$. Cells: \texttt{----+-}, \texttt{---++-}, \texttt{--+-+-}, \texttt{--+++-}, \texttt{+---+-}, \texttt{+--++-}, \texttt{+-+-+-}, \texttt{+-+++-}. Nonempty subsets by threshold: $L^\ast=0$: 72, $L^\ast=1$: 182, $L^\ast=5$: 1.
\end{description}

\medskip
The maximal sets at each level (main Table~\mref{tab:census}) are: $L=0$ the $64$ radius-one balls; $L=1,2$ the $192$ cylinders and $1{,}280$ punctured subcubes (a three-dimensional subcube minus one vertex; $V=3$ with $y=(-1,-1,-1,0,0,0)$ for the subcube $\theta_1=\theta_2=\theta_3=-1$ minus $(-1,-1,-1,1,1,1)$; first-step bound and threshold $1$; each lies in one full subcube and in one boundary ten-set); $L=3,4$ the cylinders and the boundary ten-sets; $L\ge5$ all three maximal orbits.

\subsection{Census, thresholds, obstructions}
\begin{table}[H]\centering\small
\caption{Level census for $m=4,5,6$ (main Table~\mref{tab:census} with all levels).}\label{tab:S4census}
\resizebox{\textwidth}{!}{%
\begin{tabular}{rrrrlrl}
\toprule
$m$ & $L$ & survivable & maximal & sizes & admissible & power set\\
\midrule
4 & 0 & 16 & 16 & $1\times16$ & 48 & $2^{16}$\\
4 & 1 & 48 & 32 & $2\times32$ & 48 & $2^{16}$\\
5 & 0 & 1{,}616 & 32 & $6\times32$ & 1{,}776 & $2^{32}$\\
5 & 1 & 1{,}616 & 32 & $6\times32$ & 1{,}776 & $2^{32}$\\
5 & 2 & 1{,}616 & 32 & $6\times32$ & 1{,}776 & $2^{32}$\\
5 & 3 & 1{,}696 & 112 & $4\times80$, $6\times32$ & 1{,}776 & $2^{32}$\\
6 & 0 & 7{,}072 & 64 & $7\times64$ & 853{,}680 & $2^{64}$\\
6 & 1 & 631{,}440 & 1{,}472 & $7\times1{,}280$, $12\times192$ & 853{,}680 & $2^{64}$\\
6 & 2 & 631{,}440 & 1{,}472 & $7\times1{,}280$, $12\times192$ & 853{,}680 & $2^{64}$\\
6 & 3 & 774{,}800 & 1{,}472 & $10\times1{,}280$, $12\times192$ & 853{,}680 & $2^{64}$\\
6 & 4 & 774{,}800 & 1{,}472 & $10\times1{,}280$, $12\times192$ & 853{,}680 & $2^{64}$\\
6 & 5 & 774{,}960 & 1{,}632 & $8\times160$, $10\times1{,}280$, $12\times192$ & 853{,}680 & $2^{64}$\\
\bottomrule
\end{tabular}}

\end{table}
\begin{table}[H]\centering\small
\caption{Threshold histograms over all survivable sets, number of sets with $V=\tfrac52$ exactly, and number of thresholds left uncertified by the two-sided scheme.}\label{tab:S4thr}
\resizebox{\textwidth}{!}{%
\begin{tabular}{rrlrr}
\toprule
$m$ & survivable sets & threshold histogram & $V=\tfrac52$ exactly & uncertified\\
\midrule
4 & 48 & $L^\ast=0$: 16, $L^\ast=1$: 32 & 0 & 0\\
5 & 1{,}696 & $L^\ast=0$: 1{,}616, $L^\ast=3$: 80 & 0 & 0\\
6 & 774{,}960 & $L^\ast=0$: 7{,}072, $L^\ast=1$: 624{,}368, $L^\ast=3$: 143{,}360, $L^\ast=5$: 160 & 143{,}360 & 0\\
\bottomrule
\end{tabular}}

\end{table}
\begin{table}[H]\centering\small
\caption{Obstructions: far pairs by distance with their values, admissible non-survivable sets, minimal obstructions inside the admissible family by size and value, and the largest value of any minimal obstruction (main Corollary~\mref{cor:uniform}).}\label{tab:S4obs}
\resizebox{\textwidth}{!}{%
\begin{tabular}{rp{4.2cm}rp{4.5cm}r}
\toprule
$m$ & far pairs by distance & admissible non-survivable & minimal obstructions inside the admissible family & $\max V$\\
\midrule
4 & $h=2$: 48 ($V=2$), $h=3$: 32 ($V=1$), $h=4$: 8 ($V=0$) & 0 & none & $2$\\
5 & $h=3$: 160 ($V=2$), $h=4$: 80 ($V=1$), $h=5$: 16 ($V=0$) & 80 & 4 cells: 80 ($V=2$) & $2$\\
6 & $h=4$: 480 ($V=2$), $h=5$: 192 ($V=1$), $h=6$: 32 ($V=0$) & 78{,}720 & 4 cells: 1{,}440 ($V=2$), 5 cells: 960 ($V=2$) & $2$\\
\bottomrule
\end{tabular}}

\end{table}
Representative minimal obstructions of the six-cube are the four-cell and five-cell sets
\begin{gather*}
\{(1,1,-1,1,-1,1),\ (1,1,-1,1,1,-1),\ (1,1,1,-1,-1,-1),\ (1,1,1,-1,1,1)\},\\
\{(1,1,-1,-1,1,1),\ (1,1,-1,1,-1,1),\ (1,1,-1,1,1,-1),\ (1,1,1,-1,-1,-1),\ (1,1,1,1,1,1)\},
\end{gather*}
both with $V=2$; for the four-set the uniform certificate pools to $(1,1,0,0,0,0)$, for the five-set the certificate $q=(\tfr16,\tfr16,\tfr16,\tfr13,\tfr16)$ (in the order listed) pools to $(1,1,0,0,0,0)$ as well.
